\documentclass[12pt, twoside, openright]{report} %fuente a 12pt, formato doble pagina y chapter a la derecha
\raggedbottom % No ajustar el contenido con un salto de pagina

% MÁRGENES: 2,5 cm sup. e inf.; 3 cm izdo. y dcho.
\usepackage[
a4paper,
vmargin=2.5cm,
hmargin=3cm
]{geometry}

% INTERLINEADO: Estrecho (6 ptos./interlineado 1,15) o Moderado (6 ptos./interlineado 1,5)
\renewcommand{\baselinestretch}{1.15}
\parskip=6pt

% DEFINICIÓN DE COLORES para portada y listados de código
\usepackage[table]{xcolor}
\definecolor{azulUC3M}{RGB}{0,0,102}
\definecolor{gray97}{gray}{.97}
\definecolor{gray75}{gray}{.75}
\definecolor{gray45}{gray}{.45}

% Soporte para GENERAR PDF/A
\usepackage[a-1b]{pdfx}

% ENLACES
\usepackage{hyperref}
\hypersetup{colorlinks=true,
	linkcolor=black, % enlaces a partes del documento (p.e. índice) en color negro
	urlcolor=blue} % enlaces a recursos fuera del documento en azul

% Añadir pdfs como partes del documento
\usepackage{pdfpages}

% Quitar la indentación de principio de los parrafos
\setlength{\parindent}{0em}

% EXPRESIONES MATEMATICAS
\usepackage{amsmath,amssymb,amsfonts,amsthm}

\usepackage{txfonts} 
\usepackage[T1]{fontenc}
\usepackage[utf8]{inputenc}

% Insertar graficas y fotos
\usepackage{tikz}
\usepackage{pgfplots}

\usepackage[spanish, es-tabla]{babel} 
\usepackage[babel, spanish=spanish]{csquotes}
\AtBeginEnvironment{quote}{\small}

% diseño de PIE DE PÁGINA
\usepackage{fancyhdr}
\pagestyle{fancy}
\fancyhf{}
\renewcommand{\headrulewidth}{0pt}
\fancyfoot[LE,RO]{\thepage}
\fancypagestyle{plain}{\pagestyle{fancy}}

% DISEÑO DE LOS TÍTULOS de las partes del trabajo (capítulos y epígrafes o subcapítulos)
\usepackage{titlesec}
\usepackage{titletoc}
\titleformat{\chapter}[block]
{\large\bfseries\filcenter}
{\thechapter.}
{5pt}
{\MakeUppercase}
{}
\titlespacing{\chapter}{0pt}{0pt}{*3}
\titlecontents{chapter}
[0pt]                                               
{}
{\contentsmargin{0pt}\thecontentslabel.\enspace\uppercase}
{\contentsmargin{0pt}\uppercase}                        
{\titlerule*[.7pc]{.}\contentspage}                 

\titleformat{\section}
{\bfseries}
{\thesection.}
{5pt}
{}
\titlecontents{section}
[5pt]                                               
{}
{\contentsmargin{0pt}\thecontentslabel.\enspace}
{\contentsmargin{0pt}}
{\titlerule*[.7pc]{.}\contentspage}

\titleformat{\subsection}
{\normalsize\bfseries}
{\thesubsection.}
{5pt}
{}
\titlecontents{subsection}
[10pt]                                               
{}
{\contentsmargin{0pt}                          
	\thecontentslabel.\enspace}
{\contentsmargin{0pt}}                        
{\titlerule*[.7pc]{.}\contentspage}  


% DISEÑO DE TABLAS.
\usepackage{multirow} % permite combinar celdas 
\usepackage{caption} % para personalizar el título de tablas y figuras
\usepackage{floatrow} % utilizamos este paquete y sus macros \ttabbox y \ffigbox para alinear los nombres de tablas y figuras de acuerdo con el estilo definido. Para su uso ver archivo de ejemplo 
\usepackage{array} % con este paquete podemos definir en la siguiente línea un nuevo tipo de columna para tablas: ancho personalizado y contenido centrado
\newcolumntype{P}[1]{>{\centering\arraybackslash}p{#1}}
\DeclareCaptionFormat{upper}{#1#2\uppercase{#3}\par}

% Diseño de tabla para ingeniería
\captionsetup[table]{
	format=hang,
	name=Tabla,
	justification=centering,
	labelsep=colon,
	width=.75\linewidth,
	labelfont=small,
	font=small,
}

% DISEÑO DE FIGURAS.
\usepackage{graphicx}
\graphicspath{{img/}} %ruta a la carpeta de imágenes

% Diseño de figuras para ingeniería
\captionsetup[figure]{
	format=hang,
	name=Fig.,
	singlelinecheck=off,
	labelsep=colon,
	labelfont=small,
	font=small		
}

% NOTAS A PIE DE PÁGINA
\usepackage{chngcntr} %para numeración continua de las notas al pie
\counterwithout{footnote}{chapter}

% LISTADOS DE CÓDIGO
% soporte y estilo para listados de código. Más información en https://es.wikibooks.org/wiki/Manual_de_LaTeX/Listados_de_código/Listados_con_listings
\usepackage{listings}

% definimos un estilo de listings
\lstdefinestyle{estilo}{ frame=Ltb,
	framerule=0pt,
	aboveskip=0.5cm,
	framextopmargin=3pt,
	framexbottommargin=3pt,
	framexleftmargin=0.4cm,
	framesep=0pt,
	rulesep=.4pt,
	backgroundcolor=\color{gray97},
	rulesepcolor=\color{black},
	%
	basicstyle=\ttfamily\footnotesize,
	keywordstyle=\bfseries,
	stringstyle=\ttfamily,
	showstringspaces = false,
	commentstyle=\color{gray45},     
	%
	numbers=left,
	numbersep=15pt,
	numberstyle=\tiny,
	numberfirstline = false,
	breaklines=true,
	xleftmargin=\parindent
}

\captionsetup[lstlisting]{font=small, labelsep=period}
% fijamos el estilo a utilizar 
\lstset{style=estilo}
\renewcommand{\lstlistingname}{\uppercase{Código}}

\pgfplotsset{compat=1.17} 
%-------------
%	DOCUMENTO
%-------------

\begin{document}
\pagenumbering{roman} % Se utilizan cifras romanas en la numeración de las páginas previas al cuerpo del trabajo
	
%----------
%	PORTADA
%----------	
\begin{titlepage}
	\begin{sffamily}
	\color{azulUC3M}
	\begin{center}
		\begin{figure}[H] %incluimos el logotipo de la Universidad
			\makebox[\textwidth][c]{\includegraphics[width=16cm]{Portada_Logo.png}}
		\end{figure}
		\vspace{2.5cm}
		\begin{Large}
			Grado en Ingeniería Informática\\			
			2019-2020\\
			\vspace{2cm}		
			\textsl{Apuntes}\\
			\bigskip
		\end{Large}
		 	{\Huge Inteligencia Artificial}\\
		 	\vspace*{0.5cm}
	 		\rule{10.5cm}{0.1mm}\\
			\vspace*{0.9cm}
			{\LARGE Jorge Rodríguez Fraile\footnote{\href{mailto:100405951@alumnos.uc3m.es}{Universidad: 100405951@alumnos.uc3m.es}  |  \href{mailto:jrf1616@gmail.com}{Personal: jrf1616@gmail.com}}}\\ 
			\vspace*{1cm}
	\end{center}
	\vfill
	\color{black}
		\includegraphics[width=4.2cm]{img/creativecommons.png}\\
		Esta obra se encuentra sujeta a la licencia Creative Commons\\ \textbf{Reconocimiento - No Comercial - Sin Obra Derivada}
	\end{sffamily}
\end{titlepage}

%----------
%	ÍNDICES
%----------	

%--
% Índice general
%-
\tableofcontents
\thispagestyle{fancy}

%----------
%	TRABAJO
%----------	

\pagenumbering{arabic} % numeración con múmeros arábigos para el resto de la publicación	


%----------
%	COMENZAR A ESCRIBIR AQUI
%----------	


\chapter{Tema 2: Sistemas de producción.}


  La aplicación de reglas genera conocimiento nuevo, que lleva a otras
  reglas.

 
  \textbf{Componentes:}
  

  \begin{itemize}
  \item \textbf{Base de Hechos:} Hechos que representan el problema.
    

    \begin{itemize}
    \item Son afirmaciones atómicas, referidas a términos constantes.
      
    \item Si algo no está definido en la BH es falso, Mundo cerrado.
      
    \item Representan hechos simples, objetos, lógica, etc.
      
    \item Se pueden añadir, eliminar y modificar hechos.
      
    \end{itemize}
  \item \textbf{Base de Reglas:} Reglas que representan el conocimiento para
    resolver el problema. SI condiciones ENTONCES acciones
    

    \begin{itemize}
    \item \textbf{Condiciones}:
      

      \begin{itemize}
      \item Solo existe la conjunción(y), si se quiere hacer un `o' se hace
        una regla para cada uno. Si se quiere hacer un if, se hacen 2
        reglas, una de if y otra de else.
        
      \item Las condiciones son sobre hechos de la BH.
        
      \item Contienen variables y constantes definidas en los hechos.
        
      \end{itemize}
    \item \textbf{Acciones}:
      

      \begin{itemize}
      \item Modifican la BH, añadiendo, modificando o borrando.
        
      \item Operaciones de E/S.
        
      \end{itemize}
    \end{itemize}
  \item \textbf{Motor de Inferencia:} Responsable de la ejecución de reglas.
    

    \begin{itemize}
    \item \textbf{Tipos de inferencia}:
      

      \begin{itemize}
      \item \textbf{Encadenamiento hacia delante(las que podemos completar
        ya):} Se van mirando las reglas y se comprueba de cuáles de
        ellas tenemos los antecedentes, y se ejecuta para ver si su
        resultado nos permite alcanzar a ejecutar una regla nueva. Se
        utiliza cuando hay pocos datos iniciales y muchas posibles
        conclusiones. Desventaja: No se centra en las metas. Y hay más
        comparaciones.
        \pagebreak
      \item \textbf{Encadenamiento hacia atrás(queremos completar y buscamos
        los hechos):} Se fija una submeta, que es obtener un determinado
        resultado de una regla. Observando sus antecedentes vemos que
        hechos tenemos, los que nos faltan se convierten en otra
        submeta, seguiremos este proceso hasta que encontremos una regla
        de la que tenemos todo y podamos alcanzar a obtener el resultado
        de la primera submeta, volviendo hacia atrás. Se utiliza cuando
        hay muchos datos, pero pocos relevantes. Desventaja: Hay que
        gestionar todas la metas y submetas, y no conocemos el proceso
        correcto hasta el final.
        
      \end{itemize}
    \item \textbf{Fases}:
      

      \begin{itemize}
      \item \textbf{Reducción (opcional)}
        
      \item \textbf{Equiparación:} Observar las reglas de las que se cumple
        el antecedente, con la BH actual, y estas reglas se añaden al
        Conjunto Conflicto. Puede una misma regla estar varias veces en
        el CC, si contiene distintos valores. CC=\{R1, R2(Mario),
        R2(José)\}
        
      \item \textbf{Resolución del conjunto conflicto:} Consiste en elegir
        la estrategia de ejecución de las reglas del Conjunto Conflicto.
        Para evitar bucles, una instancia se ejecuta una sola vez, se
        aplica el Principio de refracción, excepto que haya salido y
        vuelto a entrar.
        

        \begin{itemize}
        \item Estrategia por \textbf{profundidad}(la más nueva, LIFO): Las
          últimas en entrar al conjunto son las primeras en ejecutarse.
          
        \item Estrategia por amplitud(la más antigua, FIFO): Las primeras en
          entrar son las primeras en ser ejecutadas.
          
        \end{itemize}
      \item \textbf{Ejecución:} Se ejecutan las reglas del Conjunto
        Conflicto según la estrategia de resolución, y los resultados se
        van añadiendo a la Base de Hechos.
        

        \begin{itemize}
        \item Añadir: Assert o + Eliminar\textbf{:} Retract o - Modificar:
          Modify
          
        \end{itemize}
      \end{itemize}
    \item El ciclo se repite tantas veces como para vaciar el Conjunto
      Conflicto o hasta una acción de parar.
      
    \end{itemize}
  \end{itemize}

 
  \textbf{Herramientas:} 
  
  Prolog (resultado:-condiciones) y CLIPS (Mirar
  estructura en diapositivas).
  

 \pagebreak
  \textbf{Resumen representación en el espacio de estados:}
  

  \begin{itemize}
  \item Identificar el estado inicial y los posibles estados.
    
  \item Identificar los operadores para pasar de un estado a otro.
    
  \item Ver conexiones entre los estados.
    
  \item Identificar los estados finales, para saber cuando parar.
    
  \item Identificar el número de acciones ejecutadas, el coste.
    
  \item Identificar las posibles soluciones desde un estado inicial a un
    estado final, que resuelva el problema.
    
  \item Puede haber prioridad entre las reglas.
    
  \item Hay que tratar de hacer los operadores lo más genéricos posible para
    que sean flexibles.
    
  \item Describir solo las que sean necesarias, pero no comprimidas, pueden
    estar desglosadas.
    
  \end{itemize}


\chapter{Tema 3: Búsqueda}


  Son aquello problemas en los que partido de un estado inicial, por el
  que comenzaremos, y debemos llegar mediante una seria de operaciones a
  un determinado estado, el estado final. Es importante antes de
  comenzar un problema de búsqueda conocer cuáles son los estados
  posibles, y entre ellos debe estar el estado inicial y el final. La
  transición entre estos estados se realiza mediante una serie de
  operaciones. Se pueden representar como un grafo, en el que la raíz es
  el estado inicial y de el cuelgan todos los estados resultado de
  aplicar las operaciones. El número de estados posibles aumenta
  exponencialmente cuando varía el número de datos.

  \textbf{Lista abierta}: Lista de todos los nodos que aún no hemos
  explorado.

  \textbf{Parámetros importantes}:

  \begin{itemize}
  \item \textbf{Factor de ramificación, b:} Numero de posibles nodos que
    pueden derivar de otro, posibles sucesores de un nodo. El factor de
    ramificación medio es la suma de todos los factores de todos lo
    estados divididos del número de estados. Ojo tener en cuenta
    retroceso en el árbol.
    
  \item \textbf{Profundidad del árbol de búsqueda, d}: Cuantos operadores
    tenemos que aplicar para llegar a la solución, solo se puede
    calcular tras encontrarla, pero no antes.
    
  \end{itemize}

  \textbf{Búsqueda no informada}: No tenemos información de donde esta
  la solución, solo sabemos cuál es, por lo que tenemos que seguir un
  proceso mecánico lógico que recorra todos los estados. No tienen pesos
  las ramas, por lo que elegir una u otra es igual si llega a la
  solución.

  \begin{itemize}
  \item \textbf{Búsqueda en amplitud}: Sigue el orden de un FIFO, el primero
    que entra es el primero en salir se ejecutan en orden de entrada,
    una cola. Si un estado vuelve a aparecer no lo ponemos, porque ya
    esta contemplado. Este método llega al estado final por medio de
    menos operadores y eso es mejor, es la solución óptima. Como no hay
    pesos en las ramas cuando aparece la solución en una rama hemos
    acabado.
    

    \begin{itemize}
    \item Completitud(hay solución y el factor de ramificación es finito en
      cada nodo), admisibilidad (si todos los nodos tienen el mismo
      coste, encuentra la solución óptima), eficiencia y consumo de
      memoria exponencial.
      
    \end{itemize}
\pagebreak
  \item \textbf{Búsqueda en profundidad:} Sigue el orden de un LIFO. El
    último en entrar será el primero en ejecutar, una pila. Puede
    alcanzar la solución en menor tiempo, pero posiblemente lo hará
    mediante más operaciones. Se puede poner un límite de profundidad de
    las ramas, para que no se llegue a la solución con tantas
    operaciones. No se ponen estados que ya hayan aparecido. Recorre una
    rama hasta que no puede avanzar más, entonces retrocede hasta que
    pueda cansar por una rama.
    

    \begin{itemize}
    \item Requiere backtracking(cuando retrocede al llegar al límite de
      profundidad, no encuentra soluciones en esa rama o solo aparecen
      duplicados), completitud (no asegura encontrar la solución, aunque
      la haya), admisibilidad (no asegura que sea la solución óptima) y
      eficiencia(cuando la meta está lejos del estado inicial o cuando
      hay problemas de memoria).
      
    \end{itemize}
  \item \textbf{Complejidad de amplitud y profundidad:}
    

    \begin{itemize}
    \item Amplitud: Temporal y espacial exponencial $O(b^d)$
      
    \item Profundidad: Temporal exponencial $O(b^d)$ y espacial lineal
      O(b).
      
    \end{itemize}
  \item \textbf{Búsqueda de coste no uniforme}: Cuando aparecen costes en
    las transiciones. Se buscará el camino a la solución más barato en
    costes.
    

    \begin{itemize}
    \item \textbf{Dijkstra}: Amplitud pero ordenado ascendente. Consiste en
      ir escogiendo aquellas ramas que menor coste tengas, se ordena la
      lista abierta de menor a mayor y FIFO, pero si hay un estado
      repetido hay que conspirar su coste para ver si es más barato.
      Cuando vamos avanzamos hay que ir arrastrando el coste acumulado
      de las ramas ya recorridas. Es como recorrer en profundidad pero
      escogiendo el más barato. Termina cuando hemos comprobado que la
      solución encontrada es la más óptima.
      
    \item \textbf{Ramificación y acotación:} En profundidad, seguimos una
      rama cogiendo el que menos coste hasta una meta, y termina cuando
      no encontramos otra rama con menor coste que la escogida, ponemos
      el coste de esa como límite de coste de la rama, ya que si lo
      supera será mejor la anterior, pero si la encuentra esa es la
      solución.
      
    \end{itemize}

	\pagebreak
  \item \textbf{Pasos a seguir}:
    

    \begin{itemize}
    \item Formalizar el problema:
      

      \begin{itemize}
      \item Definir los posibles estados, además cuál será el inicial y el
        final o meta.
        
      \item Operaciones que se pueden realizar en los estados.
        
      \end{itemize}
    \item Estimar la complejidad, el número de estados diferentes del
      problema.
      
    \item Algoritmo de búsqueda:
      

      \begin{itemize}
      \item Orden al generar nodos.
        
      \item Lista abierta, nodos generados y nodos expandidos.
        
      \item Conocer las propiedades del algoritmo: Completitud,
        admisibilidad, complejidad temporal y espacial.
        
      \end{itemize}
    \end{itemize}
  \end{itemize}

  \textbf{Búsqueda heurística:} Se tiene conocimiento parcial sobre un
  problema que nos permite resolverlo de una manera más o menos
  eficiente. Si se tiene el conocimiento perfecto, se puede desarrollar
  un algoritmo exacto, pero si no se tiene conocimiento se hará una
  búsqueda no informada.

  \begin{itemize}
  \item Hay que hallar una función heurística h(n), que permita hallar el
    coste estimado desde el nodo n hasta un nodo objetivo, y esta
    devuelve un valor real positivo. Se descubre resolviendo modelos
    simplificados del problema real, simplificando mecánicas como
    restricciones, pero no el propio problema. Es el coste óptimo del
    problema relajado. Ejem: Distancias de Manhattan, para problemas con
    un tablero con coordenadas,
    $h(n)=| x-x_i|+| y-y_i|$, la suma de
    distancias.
    
  \item \textbf{Búsqueda Escalada (Hill Climbing)}: Básicamente buscando el
    camino que genera menores h, solo fijándonos en la h y sin acumular,
    hasta que no se pueda encontrar uno más bajo que el anterior, o
    encuentre el final. Consiste en coger un nodo de la lista abierta y
    obtener sus sucesores, calculamos la función heurística de cada uno,
    y nos quedamos en la lista abierta solo con aquel que tiene menor
    coste(no se acumulan valore, es el puro coste heurístico) y ese
    coste además debe ser menor que el del padre. No hay backtracking,
    por lo que si ese único nodo que usamos no puede avanzar, hemos
    terminado. No se analizan los nodos repetidos, ya fueron
    descartados. En cuanto un nodo sucesor sea final, el proceso
    termina.
    

    \begin{itemize}
    \item No completo (NO, ya que no siempre encuentra la solución), no
      admisible(NO, si no encuentra siempre la solución no será
      admisible) y eficiencia (es rápido y útil)
      
    \item Posibles soluciones: Poder hacer backtracking, avanzar por un par
      de ramas a la vez, reinicio aleatorio o que siga avanzando aunque
      el sucesor tenga un mayor coste.
      
    \end{itemize}
  \item \textbf{Búsqueda Mejor Primero}: Básicamente ir calculando para cada
    sucesor la suma del h propio y el coste en llegar (OJO, pero sin
    tener en cuenta los h anteriores, es la suma de transiciones pura),
    e ir avanzando por el que menor suma tiene. Consiste en ir cogiendo
    de la lista abierta el nodo de menor coste y obtener sus sucesores,
    de los que calculamos su coste que es la suma de la función
    heurística de ese nodo y el coste desde el inicio hasta llegar a ese
    nodo, y ordenamos la lista abierta en orden ascendente en coste,
    pero no desechamos ningún nodo. Este algoritmo siempre encuentra la
    solución y es la más óptima. Los nodos recorridos pasan a la lista
    cerrada, que servirá para ver si un nodo repetido lo hemos
    encontrado con menor coste, por lo que se convertiría en mejor
    camino. Si el nodo final es escogido termina, pero no solo si
    aparece como un sucesor. $f(n)=g(n)+h(n)$ El coste total es la suma
    del coste desde el inicio al nodo n y el coste desde n hasta el nodo
    final.
    

    \begin{itemize}
    \item Al final se pueden calcular los valores reales, que se indican con
      un *.
      
    \item Completitud (si existe solución la encuentra), admisibilidad
      (encuentra la solución óptima, si: sucesores finitos, los costes
      son mayores que 1 y $h(n)$ es admisible, menor que el $h^*(n)$) y
      complejidad exponencial.
      
    \item Cuando mayor $h(n)$ mejor informada esta la función heurística y
      menor numero de nodos sucesores aparecen.
      
    \end{itemize}
  \end{itemize}


\chapter{Tema 4: Redes Bayesianas}


  \textbf{Razonamiento Probabilístico:}

  \textbf{Incertidumbre}: Confianza que tenemos de un suceso.

  \begin{itemize}
  \item Se expresa como una probabilidad, que es una medida de:
    

    \begin{itemize}
    \item Proporción de veces en que algo es cierto.
      
    \item Grado de creencia en que algo es cierto.
      
    \end{itemize}
  \item \textbf{Variable aleatoria}: Que pueden tomar valores en un dominio,
    el valor asociado es desconocido, pero podemos saber la probabilidad
    de cada valor posible.
    
  \item \textbf{Espacio muestral $\Omega$:} Todos los posibles resultados de un
    experimento aleatorio.
    
  \item \textbf{Evento}: Subconjunto del espacio muestral, hay varios.
    
  \item \textbf{Evento atómico:} Evento de un único elemento y define un
    único estado.
    
  \item \textbf{Distribución de probabilidad}: Asigna a cada evento un valor
    que representa la probabilidad de que ocurra ese evento. Para hallar
    la probabilidad debemos tener datos de e. Se puede expresar como una
    vector formado por todas las probabilidades. $P(e)$
    

    \begin{itemize}
    \item Toma valores entre 0 y 1, incluidos.
      
    \item La suma de todas las probabilidades del evento es 1.
      
    \end{itemize}
  \item \textbf{Probabilidad de un evento no atómico:} Suma de las
    probabilidades de todos los eventos atómicos que lo componen, se
    suman los posibles.
    

    \begin{itemize}
    \item \textbf{$P(A)=\sum\limits P(e)$}
      
    \end{itemize}
  \item \textbf{Teorema de la probabilidad total}: La suma de las
    probabilidades que tiene de que ocurra A condicionando con otro
    evento, por la probabilidad de ese evento. \textbf{$P(A)=\sum\limits P(A, B1,B2,...)$}
    
  \item \textbf{Regla del producto:} $P(A/B)= P(A, B)/P(B)=P(A \cap B)/P(B)$
    \pagebreak
  \item \textbf{Probabilidad a posteriori(condicional)}: Sobre una
    observación, tenemos evidencias. Esa evidencia modifica el
    conocimiento sobre el dominio.
    

    \begin{itemize}
    \item \textbf{$P(A| B)=P(A \cap B)/P(B)=\alpha P(B \cap A)$}
      
    \item \textbf{Regla de la cadena: $P(A \cap B)=P(A| B)P(B)$}
      
    \item \textbf{Comparar probabilidades condicionadas:} Cuando la B es la
      misma, pero la A cambia, la suma de ambas probabilidades debe dar
      1, por lo que dejamos como incógnita $\boldsymbol{\alpha=1/P(B)}$ que es la
      \textbf{constante de normalización}. El que mayor x tenga es más
      probable siendo x el término que acompaña a $\alpha$.
      
    \item $\alpha$ permite también hallar $P(B)$ cuando no se nos da y conocemos las
      condicionadas.
      
    \end{itemize}
  \item \textbf{Probabilidad a priori:} Probabilidad que algo ocurra cuando
    no tenemos información.
    
  \item \textbf{Probabilidad conjunta:} Probabilidad de un conjunto de
    variables.
    

    \begin{itemize}
    \item Se puede representar de manera tabular, en una tabla, pero en
      problemas reales no es viable, ya que hay ciento o miles de
      variables. Hay que tener en cuenta el coste y tiempo de respuesta.
      
    \end{itemize}
  \item \textbf{Teorema de Bayes:} Se pueden calcular unas condicionales a
    través de otras.
    

    \begin{itemize}
    \item \textbf{$P(A/B)=(P(B/A)P(A))/P(B)=\alpha P(B/A)P(A)$}
      
    \end{itemize}
  \item \textbf{Independencia}: A y B lo son si y solo si la ocurrencia de
    uno de ellos no afecta a la ocurrencia del otro.
    

    \begin{itemize}
    \item $\boldsymbol{P(A| B)=P(A)}$ La probabilidad de A no se ve
      afectada por B.
      
    \item $\boldsymbol{P(A, B)=P(A)P(B)}$
      
    \item \textbf{Reducción de tamaño de distribución:} Eso implica
      algoritmo más eficiente y menos datos a especificar.
      
    \end{itemize}
  \item \textbf{Inferencia}: Calcular la probabilidad de eventos
    (probabilidad a posteriori) dada cierta evidencia. Se usa para:
    

    \begin{itemize}
    \item \textbf{Predicción.}
      
    \item \textbf{Diagnosis.}
      
    \item \textbf{Clasificación.}
      
    \item \textbf{Toma de decisiones:} Elegir acciones más útiles.
      
    \end{itemize}
  \end{itemize}

 \pagebreak
  \textbf{Redes Bayesianas:}
  

  \begin{itemize}
  \item \textbf{Independencia condicional:} Cuando hay relaciones entre
    varias variables y están relacionadas, pero por medio de otra, que
    no es causa directa. Hay un intermediario, por lo que conociéndolo
    depende del intermediario, no del que está alejado un nodo.
    
  \end{itemize}

  
$P(X,Y| Z)=P(Y| Z)P(X| Z)$ y también $P(X| Y,Z)=P(X| Z)$

Nos permite reducir el número de parámetros, normalmente de exp. a
lineal.

En las redes bayesianas se asume que hay variables ind.
condicionalmente.



  \textbf{Red Bayesianas:} Es un Grafo Acíclico Dirigido (DAG), en la
  que los nodos son las variables, los arcos indican causa directa. Cada
  nodo tiene una tabla de probabilidad condicional (CPT), indica la
  influencia de los nodos padre sobre el. Cuando no tiene es la
  probabilidad a priori.

  
  \begin{itemize}
  \item Se puede expresar de forma compacta como la probabilidad de todos
    separados por comas, distribución de probabilidad conjunta, gracias
    a que se asume independencia condicional y después podemos
    factorizarla.
    
  \item Es ir sacando términos, así:
    

    \begin{itemize}
    \item $P(A, B, C)= P(A| B,C)P(B,C)=...$
     $$P(X_1, ..., X_n)=\prod^n_{i=1} P(X_i/Padres(X_i))$$
      
    \end{itemize}
  \item La CPT tendrá $2^a$ filas en un nodo con a padres. Se pone la
    probabilidad de que sea true, el propio nodo, ya que de que sea
    false es la inversa 1-P()
    
  \item Complejidad $O(n\cdot 2^b)$ n variables y b padres.
    
  \end{itemize}

 
  \textbf{Causalidad vs. Correlación:} Que numéricamente estén
  relacionados, no quiere decir que existan una relación directa, pueden
  existir variables intermedias. Fijarse si esto ocurre.
  

  \begin{itemize}
  \item Correlación no implica causalidad, pero Causalidad implica
    correlación.
    
  \end{itemize}

 
  \textbf{Inferencia}: Calcular la probabilidad a posteriori de una
  variable dada cierta evidencia.
  

  \begin{itemize}
  \item Hay que \textbf{tener en cuenta}: La variable pregunta, las
    variables evidencia y las ocultas.
    \pagebreak
  \item \textbf{Inferencia exacta por enumeración}:
    

    \begin{itemize}
    \item 1. Se aplica la regla del producto. Donde $\alpha=\frac{1}{\textit{lo de la derecha}}$ o
      $\frac{1}{\sum\limits_s \alpha P_s} $
      

      \begin{itemize}
      \item \textbf{$P(A| B)=\alpha P(A, B)$.}
        
      \end{itemize}
    \item 2. Después aplicamos la regla de la probabilidad total, OJO con
      los sumatorios y poder sacar variables para simplificar o que sume
      1. En general las variables que no son ancestros de variables
      pregunta o variables evidencia son irrelevantes.
      

      \begin{itemize}
      \item \textbf{$\alpha P(A,B)=\alpha \sum\limits P(A, B, C)$.}
        
      \end{itemize}
    \item 3. Lo de dentro del sumatorio se factoriza y operamos.
      

      \begin{itemize}
      \item Las evidencias tendrá el valor fijo, pero las ocultas se
        contemplan con el sumatorio. Añade en forma de sumatorio las
        variables ocultas. Cuando no se indica si la variable pregunta
        es true o false, se hallan ambas posibilidades por separado,
        esto a su vez nos permite hallar $\alpha$, suman 1 ambas prob.
        
      \end{itemize}
    \item El problema de este método son la variable ocultas que tienen
      complejidad exponencial.
      
    \item Es eficiente para poliarboles, que como mucho haya un arco entre
      cada par de nodos. Esta estructura da lugar a una complejidad
      lineal.
      
    \item No es eficiente cuando la estructura de la red no es un poliarbol,
      pasa de ser complejidad lineal a ser exponencial. En estos casos
      se usa la inferencia aproximada.
      
    \end{itemize}
  \item \textbf{Inferencia aproximada}: Se hace mediante estimaciones, para
    ello se hace un muestreo de la red Bayesiana, comenzando desde
    padres hasta los hijos. Muestrear, quiere decir que siguiendo la
    distribución de probabilidad sacamos un número aleatorio. Por
    ejemplo $P(Nublado)=(0.5, 0.5)$ Sacamos un numero entre 0 y 1, si cae
    en la parte de la izquierda es true, si no false. $0-0.5$ true y $0.5-1$
    false.
    

    \begin{itemize}
    \item De esta manera sacamos que valor tiene cada variable y podemos
      centrarnos en ese caso concreto, y hallar su probabilidad como
      haríamos en el caso anterior.
      
    \item \textbf{Muestreo directo:} La probabilidad de un evento se estima
      como el numero de casos del evento generados por muestreo dividido
      entre el numero total de casos muéstrales.
      
	  $$P(x_1, x_2, ..., x_n) \simeq \frac{\#casos(x_1, x_2, ..., x_n))}{\#total Casos}$$
	  $$P(X/e) = \frac{\#casos(X, e))}{\#casos(e)}$$

    \end{itemize}
  \end{itemize}

 
  \textbf{Clasificador \emph{Naïve Bayes}}: Caso especial de red
  bayesiana con una estructura en la que de un nodo cuelgan el resto. El
  nodo raíz es del que deseo conocer la probabilidad, y los hijos son
  los atributos y son independientes entre ellos.
  $$p(c_i/a_1 ... a_n)= \frac{P(a_1 ... a_n / c_i)p(c_i)}{p(a_1 ... a_n)} \rightarrow p(c_i/a_1 ... a_n)= \alpha p(C= c_i) \prod_k P(A_k = a_k / C= c_i)$$
  \begin{itemize}
  \item \textbf{Se puede resolver con} Bayes normal o mediante los métodos
    de inferencia.
    
  \end{itemize}

  


  \textbf{Razonamiento Probabilístico en el Tiempo:}

  
  \begin{itemize}
  \item Hasta el momento hemos contemplado solo mundos estaticos, pero ahora
    consideramos el paso del tiempo. Cuando el mundo es no determinista,
    que hay varias posibles decisiones en cada paso, se emplea el
    procesos de decisión de Markov.
    
  \item \textbf{Modelos de Markov:} Se representa con un grafo la
    probabilidad de las transiciones, estando en un tiempo t y yendo a
    t+1, el momento siguiente. El estado actual determina la
    distribución de probabilidad del estado siguiente. El momento
    inicial será $E_0$.
    

    \begin{itemize}
    \item \textbf{Hipótesis de Markov:} Se asume que el estado actual solo
      depende del anterior, es condicionalmente independiente de los que
      no son el inmediato anterior.

	  $$P(E_{t+1}=s1 / E_t =s2)$$
      
    \item Se puede representar como una red bayesiana, en la que los estado
      que no dependen de otros tienen su probabilidad a priori y las que
      dependen de otro su tabla de probabilidad condicionada.
      
    \item \textbf{Calcular la probabilidad:} Se hace como inferencia exacta.
      Lo primero es hacer la regla del producto, después hacer la
      probabilidad total con sumandos y finalmente factorizar. Pero esta
      opción es muy pesada ya que hay que considerar todos los posibles
      caminos anteriores.

	  $$P(E_t = s_j)= \sum\limits_{\textit{Todos los caminos que acaban en el estado} P(E_0, E_1, ..., E_t = s_j)}$$
      
    \item \textbf{Programación dinámica:} Es una técnica/truco, que consiste
      en hacer una definición recursiva en función del instante
      anterior. Se comienza desde el final, y se van llamando unas a
      otras hasta el inicial. Algoritmo de Simulación hacia delante, que
      tiene complejidad lineal para un tiempo t.

	  $$\forall j, P(E_{t+1}=s_j) = \sum\limits^N_{i=1} P(E_{t+1}=s_j/E_t=s_i)P(E_t = s_i)$$
      
    \item \textbf{Distribuciones estacionarias:} Normalmente solo podemos
      predecir a corto plazo, a medida que nos alejamos de los datos
      iniciales menos sabemos, llegando en el infinito a ser
      equiprobable y no nos permite decantarnos por ninguna.
      

      \begin{itemize}
      \item La incertidumbre se acumula, hasta que no sabremos cual es el
        estado.
        
      \item Para la mayoría de las cadenas la distribución al final es
        independiente de la inicial, la distribución al final es igual
        si empezamos con $P(X_0=lluvia)=1$
        
      \end{itemize}
    \end{itemize}
  \item \textbf{Modelos de Markov Ocultos (HMMs):} Hay variables
    observables, evidencias, que determinan el estado. Tiene forma de
    red bayesiana en el que las observaciones cuelgan del estado y solo
    dependen del estado actual. Cada estado depende del inmediato
    anterior.
    

    \begin{itemize}
    \item \textbf{Se asume:}
      

      \begin{itemize}
      \item Es un proceso estacionario.
        
      \item Se cumple la hipotesis de Markov.
        
      \item Las observaciones en t solo dependen del estado en t. P(O/E)
        
      \end{itemize}
    \item \textbf{Se define por:}
      

      \begin{itemize}
      \item Conjunto de estados y de observaciones.
        
      \item La probabilidad a priori, aquel que no depende del anterior.
        $P(E_0)$
        
      \item El modelo de transiciones. $P(E_1/E_0)$
        
      \item El modelo de observaciones. $P(O_1/E_1)$
        
      \end{itemize}
    \item La probabilidad de las transiciones se representan en la tabla de
      probabilidad condicionada.
      
    \item \textbf{Inferencia en HMMs:}
      

      \begin{itemize}
      \item Se puede resolver por definición recursiva.
        
      \item La forma mas eficiente es por programación dinámica, pero no la
        veremos.
        
      \item Nosotros las resolveremos mediante inferencia exacta en redes
        Bayesianas, a pesar de que es menos eficiente.
        
      \item \textbf{Se factoriza como:} 
	  
	  $P(E_0, ..., E_t, O_0, ..., O_t) = P(O_0/E_0)P(E_0) \prod^T_{t=1} P(E_t/ E_{t-1})P(O_t/E_t)$
        
      \end{itemize}
    \end{itemize}

	\pagebreak
  \item Tareas típicas de Inferencia:
    

    \begin{itemize}
    \item \textbf{Problema de Evaluación}: Calcular la probabilidad de una
      secuencia de observaciones.
      
    \item \textbf{Problema de Decodificacion}: Dada una secuencia de
      observaciones, determinar cual es la secuencia de estados
      correspondiente que explica mejor esas observaciones.
      
    \item \textbf{Problema de Filtrado:} Distribución de probabilidad del
      estado actual dada cierta evidencia histórica, desde t hacia
      atrás.
      
    \item \textbf{Problema de Predicción:} Probabilidad de estados futuros
      dada evidencia.
      
    \end{itemize}
  \item \textbf{Procesos de Decisiones de Markov (MDPs)}:
    

    \begin{itemize}
    \item Se toma una decision en cada instante. El caso general es que sean
      acciones no deterministas, hay varias posibilidades en cada caso y
      cada una de ellas tiene un probabilidad y debemos decidir cuál
      tomamos.
      
    \item Las \textbf{probabilidades de transicion} dependen de las
      acciones.
      
	  $$P(S_{t+1} = s' / S_t = s, A_t = a)= P(s'/s, a)$$
    \item Hay un \textbf{refuerzo o un coste}, indica como de bueno o malo
      es esa transición, por cada par de estado-accion:

	  $$R(S_t = s, A_t = a) = R(s, a)$$
      
    \item Se define como una tupla \textbf{\textless S, A, P,
      R\textgreater:}
      

      \begin{itemize}
      \item \textbf{S}, estados.
        
      \item \textbf{A}, acciones.
        
      \item \textbf{P}, probabilidades P(s'/ s, a). Cada estado tiene una
        probabilidad para cada accion con sus estados contiguos.
        
      \item \textbf{R}, refuerzos o costes R(s, a).
        
      \end{itemize}
    \item El objectivo es determinar que accion ejecutar en cada estado para
      maximizar el refuerzo o minimizar el coste.
      
    \item Al no ser determinista hay \textbf{dos opciones:}
      

      \begin{itemize}
      \item \textbf{Politica} (Acciones estocásticas): Se determinada para
        cada estado un accion determinada.
        

        \begin{itemize}
        \item Es un mapeo completo de estados a acciones, pero no es una
          secuencia de acciones. Aunque haya un fallo en la ejecucion el
          agente puede seguir.
          
        \item Maximiza el refuerzo esperado o minimiza el coste esperado en
          lugar de alcanzar un estado meta.
          
        \item Para cada MDP existe una politica optima.
          
        \item Determina que hacer independientemente del efecto de cualquier
          accion en cualquier instante de tiempo.
          
        \end{itemize}
      \item \textbf{Replanificar}: Se vuelve a planificar cuando sale del
        anterior, de esta manera cada estado tiene en cuenta el
        anterior.
        
      \end{itemize}
    \item \textbf{Accion estocástica}: Consigue el efecto deseado con
      probabilidad p.
      
    \item Para el caso de \textbf{minimizar coste:}
      

      \begin{itemize}
      \item \textbf{Politica optima $\pi$*}: mínimo coste esperado.
        

        \begin{itemize}
        \item \textbf{Acciones deterministas}: Si la accion a lleva al
          estado s', el coste es:
          

         
		  
            \textbf{C(a) + costeDesde(s')}
          
			
        \item \textbf{Acciones no deterministas}: Si la accion a tiene
          efectos probabilisticos, el coste esperado es:
          

          
		  
            \textbf{C(a) + $\sum\limits$sub(s') P(s'|{} s, a) *
            costeDesde(s')}
          
			
        \end{itemize}
      \end{itemize}
    \item \textbf{Funcion de valor:} coste esperado de un estado.
      

      \begin{itemize}
      \item \textbf{V(s):} coste esperado de alcanzar la meta desde s.
        
      \item \textbf{Busqueda}: Coste del camino optimo desde s. Se puede
        usar como heuristica, y es la heurística perfecta $h^*(s)$
        
      \item \textbf{MDP}: Coste esperado de la estrategia optima para
        alcanzar la meta desde s. Conociendo V(s) podemos calcular una
        politica optima $\pi^*$.
        
      \end{itemize}
    \item \textbf{Calcular V(s).}
      

      \begin{itemize}
      \item \textbf{Ecuaciones de Bellman}. Acciones deterministas (Es
        programación dinámica, recursividad)
        

        \begin{itemize}
        \item Si s es un estado meta, $V(s) = 0$. Es sumidero. Inicialmente
          todos están a 0.
          
        \item En el resto de casos, donde s' es el estado resultado de
          aplicar a en s:
          
		  $$V(s) = \min_{a\in A(s)} [c(a)+V(s')]$$

        \end{itemize}
      \item \textbf{Ecuaciones de Bellman}. Acciones no deterministas.
        

        \begin{itemize}
        \item Dominios estocasticos, apartir de la accion a: $c(a) + \sum\limits_{s'\in S} P_a (s'/s)V(s')$
          
        \item Entonces, el estado meta $V(s)=0$ para costes y en el resto de
          casos:

		  $$V(s)_{i+1}= R(s)+ \max_{a\in A(s)} [\gamma \sum\limits_{s'\in S} P_a(s'/s) V(s')]$$
          
        \item \textbf{Politica optima}: Consiste en elegir las accione que
          han generado el mínimo.
          $$\pi^*(s)= \arg \max_{a} [\gamma \sum\limits_{s'\in S} P_a(s'/s) V^*(s')]$$

		  \pagebreak
        \item \textbf{Resolviendo las ecuaciones de Bellman:}
          

          
		  
            \textbf{Algoritmo de Iteración de Valor:}

			
            
			
              El $V(s)$ de los finales es 0 para costes y el refuerzo si
              hay refuerzos.
             
			  
              Inicialmente todos se ponen a 0.
             
			  
              Se hacen tantas rondas como sean necesarias para que los
              valores de $V(s)$ se estabilicen entre rondas. Y cada ronda
              calcula $V(s)$ para cada estado.
              
			  
              Termina cuando alcanza un punto fijo, cuando los valores
              no cambian en dos iteraciones sucesivas. Se estabiliza.
            
			  
        \end{itemize}
      \end{itemize}
    \end{itemize}
  \end{itemize}

 
  \textbf{Lógica Borrosa:}
  

  \begin{itemize}
  \item Los conceptos no son ciertos o falsos de forma clara, a diferencia
    de la lógica clásica en la que los valores solo son true o false. Se
    emplean conceptos y modificadores que son difusos.
    
  \item \textbf{Representación:}
    

    \begin{itemize}
    \item Conjunto borroso.
      
    \item \textbf{Función de pertenencia:} Indica en que medida un elemento
      pertenece al conjunto borroso. Valores en el rango {[}0, 1{]}, en
      el que 0 representa absolutamente falso y 1 absolutamente
      verdadero.
      
    \end{itemize}
  \item \textbf{Lógica borrosa vs. probabilidad:}
    

    \begin{itemize}
    \item \textbf{Probabilidad}: Lo hechos ocurren o no. Expresan
      conocimiento parcial.
      
    \item \textbf{Lógica borrosa}: Conceptos vagos, inciertos. Expresa grado
      de verdad parcial.
      
    \end{itemize}
  \item \textbf{Sistema de reglas borrosas:}
    

    \begin{itemize}
    \item Descripción del problema, las reglas que aparecen.
      
    \item Definir los términos borrosos de las reglas, mediante su función
      de pertenencia.
      

      \begin{itemize}
      \item \textbf{Variables}: Expresan cualidades.
        
      \item \textbf{Valores}: Las variables toman valores de un dominio
        discreto. Los límites entre los valores son borrosos.
        
      \item \textbf{Tipos}: Sigmoide(la S), Gausiana (n), triangular y
        trapezoidal. Las dos últimas son las que utilizaremos nosotros.
        
      \item \textbf{Modificadores}: Operan sobre la función de pertenencia.
        ², ³...
        
      \end{itemize}
	  \pagebreak
    \item \textbf{Combinar términos:}
      

      \begin{itemize}
      \item \textbf{Conjunción}: Intersección. El menor de ambas para cada
        x. MIN
        
      \item \textbf{Disyunción}: Unión. El máximo en cada x de las
        funciones. MAX
        
      \item \textbf{Negación}: Los valores inversos de cada x. 1-valor
        
      \end{itemize}
    \item \textbf{Combinar reglas}, para generar una salida única. (Varias
      son parcialmente cercas)
      

      \begin{itemize}
      \item p-\textgreater q p es cierto en un grado, entonces q también es
        cierto en un grado.
        
      \end{itemize}
    \end{itemize}
  \item \textbf{Sistema de reglas borrosas:}
    

    \begin{itemize}
    \item Entrada datos nítida.
      
    \item Borrosificar el dato nítido.
      
    \item Base conocimiento, reglas borrosas.
      
    \item Desborrosificar, pasar a un valor nítido.
      
    \item Salida dato.
      
    \end{itemize}
  \item \textbf{Inferencia con reglas borrosas:} 4 pasos.
    

    \begin{itemize}
    \item \textbf{Método de Mandami}: es el método típico de inferencia
      borrosa.
      

      \begin{itemize}
      \item \textbf{Borrosificar las entradas:}
        

        \begin{itemize}
        \item Determinar en que grado la entrada nítida pertenece a los
          conjuntos borrosos, para cada valor que puede tomar la
          variable.
          
        \end{itemize}
      \item \textbf{Evaluación de reglas:}
        

        \begin{itemize}
        \item En que medida las entradas borrosificadas verifican los
          antecedentes de la regla, se evalúan todas la reglas. AND es
          intersección, OR es unión y NOT la inversa. Si es un solo
          valor se coge directamente el grado de la entrada.
          
        \item Se obtiene la \textbf{Similitud}, que es el mayor grado que
          puede obtener el consecuente.
          
        \item \textbf{El consecuente}, resultado de cortar la función de
          pertenencia del consecuente al nivel que marca la similitud
          del antecedente.
          
        \end{itemize}
      \item \textbf{Agregación de los consecuentes:}
        

        \begin{itemize}
        \item Unificación de las salidas de todas las reglas tras evaluarlas
          en un solo conjunto borroso. Nosotros hacemos el MAX para cada
          x, unión.
          
        \end{itemize}
      \item \textbf{Desborrosificar el resultado:}
        

        \begin{itemize}
        \item Convertir el resultado en un valor nítido, el más común es
          centro de gravedad o centro de del área.
          
        \end{itemize}
      \end{itemize}
    \end{itemize}
  \end{itemize}

\chapter{Tema 5: Aprendizaje Automático}
\href{https://drive.google.com/file/d/1LI9yhfP9Ln8pNsRSsPNDDnHGNqfDMPcd}{Acceso en Drive a las diapositivas}

\chapter{Tema 6: Tecnicas bioinspiradas}
\href{https://drive.google.com/file/d/1yPlKdsKnWpFO3Jqh7U2IcJ1cHM6H5vnZ}{Acceso en Drive a las diapositivas}

\chapter{Tema 7: Robótica}
\href{https://drive.google.com/file/d/1yMhsPLLKCZ4Nap7Vca12DfdT8iD_fqJf}{Acceso en Drive a las diapositivas}

\end{document}